\documentclass[conference]{IEEEtran}
% Extended Abstract metadata
\title{Development, simulation \& implementation of precoding techniques for satellite links}
\author{
Wannes Baes\\[0.2cm]
Supervisor: Prof.\ Dr.\ ir.\ Nele Noels (Ghent University)\\
Counsellor: Dr.\ ir.\ Rodrigo Moraes (Antwerp Space)
}

% Citations
\usepackage{cite}

% Math
\usepackage{amsmath,amssymb,amsfonts}

% Algorithms
\usepackage{algorithmic}

% Figures
\usepackage{graphicx}

% Symbols
\usepackage{textcomp}

% Abbreviations
\usepackage{acronym}
\acrodef{CSI}{channel state information}
\acrodef{MIMO}{multiple-input multiple-output}
\acrodef{MU}{multi-user}

\begin{document}

% Title
\maketitle

% Abstract
\begin{abstract}
Linear precoding has become a cornerstone of \acs*{MIMO} communications, as it enables spatial multiplexing gains that translate directly into higher spectral efficiency.
These benefits are particularly valuable in multi-beam high-throughput satellite systems, where aggressive frequency reuse across the narrow spot beams turns the forward link into an interference-limited \acs*{MU}-\acs*{MIMO} channel. 
In this environment, precoding is crucial for mitigating inter-beam interference.
In mobile satellite scenarios, however, the combination of long propagation delays and time-varying channels induced by user mobility introduce complications.
More specific, the \acs*{CSI} available during the precoder computation is generally outdated with respect to the channel experienced at the time of transmission.
This work investigates the performance degradation caused by this \acs*{CSI} mismatch and proposes a low-complexity compensation strategy.
This work quantifies the impact of the \acs*{CSI} feedback delay on the system performance, revealing a substantial performance loss.
To counter this degradation, a low-complexity prediction strategy is proposed.
Numerical results show that this scheme recovers a significant fraction of the performance loss under the considered channel model.
\end{abstract}

% Keywords
\begin{IEEEkeywords}
\acs*{MIMO}, Precoding Techniques, Mobile Satellite Systems, Outdated \acs*{CSI}
\end{IEEEkeywords}


\section{Introduction}
% ====================
%   Introduction
% ====================
citation~\cite{signal_processing_for_hts}

\section{System Model}
% ====================
%   System Model
% ====================


\section{Precoder Design}
% ====================
%   Precoder Design
% ====================


\section{Outdated CSI Compensation}
% ================================
%   Outdated CSI Compensation
% ================================


\section{Simulation Results}
% ==========================
%    Simulation Results
% ==========================


\section{Conclusion}
% ====================
%    Conclusion
% ====================


\bibliographystyle{IEEEtran}
\bibliography{references}

\end{document}
